% SHARD-ID: CA-24-GAUSSIAN-BOSON-SYMBOL-CALCULUS
% SHARD-TITLE: Gaussian Boson Symbol Calculus and the Boost-Time Residual
% SHARD-SUMMARY: Defines the first scalar Gaussian-boson coefficient tier and its Fourier dispersion symbol in d=1,2,3.
% SHARD-SUMMARY: Derives the one-particle boost-time commutator symbol as one half of the gradient of omega squared.
% SHARD-SUMMARY: Checks the nearest-neighbour lattice Klein-Gordon example and its small-momentum Lorentz limit in Julia.
% SHARD-KEYWORDS: Gaussian boson, Fourier symbol, dispersion, boost, Klein-Gordon, Julia check

\section{Gaussian Boson Symbol Calculus and the Boost-Time Residual}
\label{sec:gaussian-boson-symbol-calculus}

\subsection{Status, Sources, and Dependencies}

\textbf{Status: sourced model class plus checked flat local derivation.}  The
standard lattice scalar-field Hamiltonian and dispersion are sourced.  The
boost-time residual calculation is a local one-particle derivation on flat
\(L^2(dk)\), checked by Julia on positive-dispersion patches.  It is not yet a
sourced OAR or mass-shell Poincare representation theorem, and massless
zero-mode examples are coefficient-level tests only.

Dependencies: CA-11, CA-12, CA-23, and CONVENTIONS.md (h), (j), (k).

% Source: literature/md/2010.11121/2010.11121.md:598 -- free lattice Hamiltonian for scalar fields.
% Source: literature/md/2010.11121/2010.11121.md:603 -- lattice dispersion gamma_mu_N.
% Source: literature/md/2010.11121/2010.11121.md:614 -- physical mass convention mu_N^2 = epsilon_N^2 m^2 + 2d.
% Source: literature/md/2010.11121/2010.11121.md:623 -- lattice Klein-Gordon dispersion gamma_m^(N)(k)^2.
% Source: literature/md/2010.11121/2010.11121.md:648--681 -- sourced OAR one-particle scalar product, flat Fock realization, dynamics, speed 1, and translations.
% Source: literature/md/2010.11121/2010.11121.md:676 -- continuum dynamics satisfies the Klein-Gordon equation.
% Source: literature/md/2010.11121/2010.11121.md:681 -- Klein-Gordon propagation speed 1 and spatial translations.
% Source: literature/md/2010.11121/2010.11121.md:1626--1628 -- massless fields require a zero-mode policy.
% Source: references/cft/Schottenloher2008/Schottenloher2008.md:4186 -- free bosonic QFT example assumes m > 0.
% Source: references/cft/Schottenloher2008/Schottenloher2008.md:4244 -- sourced free-boson construction has a natural Poincare action on Fock space.
% Convention: CONVENTIONS.md (j) -- CA-24--CA-26 use flat L^2(dk) local symbol algebra, not yet the sourced OAR/mass-shell theorem.
% Convention: CONVENTIONS.md (k) -- self-adjoint Poincare commutator signs.
% Checked: test/runtests.jl, "Gaussian boson Klein-Gordon symbols" and "Gaussian boson boost-time symbols".

\subsection{Coefficient Tier}

Fix \(d\in\{1,2,3\}\).  On a periodic hypercubic lattice with spacing
\(\epsilon\), use canonical fields
\[
  [q_x,q_y]=0,\qquad [p_x,p_y]=0,\qquad [q_x,p_y]=i\delta_{xy}.
\]
The first scalar Gaussian tier is
\[
  H=\frac{1}{2}\sum_x p_x^2
    +\frac{1}{2}\sum_{x,r}q_xV_rq_{x+r},
  \qquad V_r=V_{-r}\in\mathbb R.
\]
Quasi-locality means that \(V_r\) has enough decay, or finite range in the first
examples, so that the Fourier symbol and its derivatives used below exist.

Define
\[
  \omega_\epsilon(k)^2
  :=
  \sum_r V_r e^{i\epsilon k\cdot r}.
\]
For a stable scalar model, this symbol is non-negative, with the zero-mode
policy stated separately in the massless case.  The helper parameter \(m\) is a
nonnegative mass label; the generator statements below use only the stricter
case \(\omega_\epsilon(k)>0\) on the chosen patch.

\subsection{Nearest-Neighbour Klein-Gordon Example}

The sourced lattice scalar-field example has, in the notation of this shard,
\[
  V_0=m^2+\frac{2d}{\epsilon^2},
  \qquad
  V_{\pm e_a}=-\frac{1}{\epsilon^2}.
\]
Therefore
\[
\begin{aligned}
  \omega_{\epsilon,m}(k)^2
  &=
    m^2+\frac{2d}{\epsilon^2}
    -\frac{1}{\epsilon^2}\sum_{a=1}^d
      (e^{i\epsilon k_a}+e^{-i\epsilon k_a})\\
  &=
    m^2+\frac{2}{\epsilon^2}\sum_{a=1}^d(1-\cos(\epsilon k_a)).
\end{aligned}
\]
For fixed physical \(k\) and \(\epsilon\to0\),
\[
  \omega_{\epsilon,m}(k)^2
  =
  m^2+|k|^2+O(\epsilon^2|k|^4).
\]

\subsection{Momentum-Space Position}

After diagonalization, this shard uses the flat local momentum space
\(L^2(\mathcal U,dk)\).  The one-particle Hamiltonian is multiplication by
\(\omega(k)\), with \(\omega(k)>0\) on \(\mathcal U\), and
\[
  X_a=i\partial_{k_a}
\]
on a common smooth core.  The time-zero boost candidate corresponding to the
energy first moment is the flat-measure symmetrized operator
\[
  K_a=\frac{1}{2}(X_a\omega+\omega X_a).
\]
The symmetrization is the one-particle version of taking a self-adjoint first
moment of the energy density.  If the one-particle scalar product is changed,
the adjoint and hence the boost formula must be rederived; no unitary
equivalence with the OAR weighted space or Schottenloher mass-shell realization
is asserted here.

\subsection{Boost-Time Residual}

For a test wavefunction \(\psi\),
\[
\begin{aligned}
  [K_a,\omega]\psi
  &=
  \frac{1}{2}(X_a\omega^2-\omega X_a\omega)\psi\\
  &=
  i\,\omega\,(\partial_a\omega)\psi.
\end{aligned}
\]
Thus
\[
  i[\omega,K_a]
  =
  \omega\,\partial_a\omega
  =
  \frac{1}{2}\partial_a\omega(k)^2.
\]
This is the algebraic output of the boost-time commutator.  With the
vector-field-to-Stone sign map of CONVENTIONS.md (k), the continuum Poincare
target is \(i[H,K_a]=P_a\).  Taking \(P_a=k_a\), one needs
\[
  \frac{1}{2}\partial_a\omega(k)^2=k_a.
\]
Integrating over \(k\) gives
\[
  \omega(k)^2=m^2+|k|^2
\]
on the connected low-energy patch.  With speed \(c\), the right side is
\(c^2k_a\), giving \(\omega(k)^2=m^2+c^2|k|^2\).

\subsection{Lattice Residual}

For the nearest-neighbour lattice Klein-Gordon symbol,
\[
  \frac{1}{2}\partial_a\omega_{\epsilon,m}(k)^2
  =
  \frac{\sin(\epsilon k_a)}{\epsilon}.
\]
Therefore the boost-time commutator has the correct continuum limit:
\[
  \frac{\sin(\epsilon k_a)}{\epsilon}=k_a+O(\epsilon^2k_a^3).
\]
The Julia tests check this in \(d=1,2,3\), both from the closed formula and by
evaluating the coefficient symbol
\(\sum_rV_re^{i\epsilon k\cdot r}\).

\subsection{Rotation Residual}

For the one-particle rotation candidate
\[
  J_{ab}=X_aP_b-X_bP_a,\qquad P_a=k_a,
\]
one obtains
\[
  i[\omega,J_{ab}]
  =
  k_b\partial_a\omega-k_a\partial_b\omega.
\]
Thus rotations commute with the Hamiltonian exactly when the dispersion is
radial on the momentum region under consideration.  For a hypercubic lattice
this is not exact over the full Brillouin zone, but the Klein-Gordon expansion
is isotropic to quadratic order at \(k=0\).

\subsection{Boundary}

This shard is a one-particle symbol calculation.  The many-body generators are
second quantizations plus vacuum-energy and domain policies.  Real-space
locality of \(K_a=\sum_xx_ah_x\), boundary terms, and convergence of the
second-quantized operators are deferred to later shards.
